
\documentclass{article}
\usepackage{colortbl}
\usepackage{makecell}
\usepackage{multirow}
\usepackage{supertabular}

\begin{document}

\newcounter{utterance}

\centering \large Interaction Transcript for game `wordle', experiment `high\_frequency\_words\_no\_clue\_no\_critic', episode 6 with mistral{-}large{-}3{-}azure.
\vspace{24pt}

{ \footnotesize  \setcounter{utterance}{1}
\setlength{\tabcolsep}{0pt}
\begin{supertabular}{c@{$\;$}|p{.15\linewidth}@{}p{.15\linewidth}p{.15\linewidth}p{.15\linewidth}p{.15\linewidth}p{.15\linewidth}}
    \# & \multicolumn{2}{c}{Player} && \multicolumn{2}{c}{Game Master} \\
    \hline

    \theutterance \stepcounter{utterance}  
    & & & \multicolumn{4}{p{0.6\linewidth}}{
        \cellcolor[rgb]{0.9,0.9,0.9}{
            \makecell[{{p{\linewidth}}}]{
                \texttt{\tiny{[P1$\langle$GM]}}
                \texttt{Du är en språkmagiker som gillar att gissa ord genom att använda de givna reglerna.} \\
\\ 
\texttt{Välkommen till Wordle! Du har sex försök på dig att gissa målordet, ett giltigt svenskt ord med fem gemena bokstäver. Använd taggarna "förklaring:" och "gissning:" för att ge en kortfattad förklaring till varje gissning.} \\
\\ 
\texttt{Till exempel, om din gissning är "fråga", ska ditt svar vara} \\
\texttt{förklaring: detta är ett vanligt svenskt ord med fem bokstäver, och jag börjar min gissning med detta ord.} \\
\texttt{gissning: fråga} \\
\\ 
\texttt{Efter varje gissning kommer ditt svar att valideras, och du kommer att få återkoppling som visar vilka bokstäver som är korrekta (gröna), vilka bokstäver som är korrekta men på fel position (gula), och vilka bokstäver som är felaktiga (röda). Denna återkoppling kan vara användbar för att avgöra vilka bokstäver du ska ta med eller utesluta i din nästa gissning.} \\
\\ 
\texttt{Till exempel kan återkopplingen för "fråga" vara:} \\
\texttt{gissning\_återkoppling: f<gul> r<gul> å<grön> g<gul> a<röd>} \\
\\ 
\texttt{Förklaringen ska innehålla detaljer om hur gissning\_återkoppling används för att komma fram till en ny gissning.} \\
\\ 
\texttt{Generera ingen annan text.} \\
\\ 
\texttt{Låt oss börja med din första gissning.} \\
            }
        }
    }
     \\ \\

    \theutterance \stepcounter{utterance}  
    & \multicolumn{4}{p{0.6\linewidth}}{
        \cellcolor[rgb]{0.561,0.737,0.561}{
            \makecell[{{p{\linewidth}}}]{
                \texttt{\tiny{[P1$\rangle$GM]}}
                \texttt{förklaring: Jag börjar med ett vanligt svenskt ord som innehåller flera vanliga vokaler och konsonanter för att maximera informationen om målordet.} \\
\texttt{gissning: solen} \\
            }
        }
    }
    & & \\ \\

    \theutterance \stepcounter{utterance}  
    & & & \multicolumn{2}{p{0.3\linewidth}}{
        \cellcolor[rgb]{0.9,0.9,0.9}{
            \makecell[{{p{\linewidth}}}]{
                \texttt{\tiny{[GM$|$GM]}}
                \texttt{attempts = 1} \\
\texttt{target\_word = panna} \\
\texttt{guess = solen} \\
\texttt{guess\_feedback = s<red> o<red> l<red> e<red> n<yellow>} \\
            }
        }
    }
    & & \\ \\

    \theutterance \stepcounter{utterance}  
    & & & \multicolumn{4}{p{0.6\linewidth}}{
        \cellcolor[rgb]{0.9,0.9,0.9}{
            \makecell[{{p{\linewidth}}}]{
                \texttt{\tiny{[P1$\langle$GM]}}
                \texttt{gissning\_återkoppling: s<röd> o<röd> l<röd> e<röd> n<gul>} \\
\\ 
\texttt{Ge ditt svar endast i detta format.} \\
\texttt{förklaring: detaljer} \\
\texttt{gissning: ord} \\
            }
        }
    }
     \\ \\

    \theutterance \stepcounter{utterance}  
    & \multicolumn{4}{p{0.6\linewidth}}{
        \cellcolor[rgb]{0.561,0.737,0.561}{
            \makecell[{{p{\linewidth}}}]{
                \texttt{\tiny{[P1$\rangle$GM]}}
                \texttt{förklaring: Bokstäverna *s, o, l, e* är helt fel och utesluts. *n* finns med i ordet men på fel position. Nästa gissning testar nya vanliga bokstäver samt inkluderar *n* på en annan plats.} \\
\texttt{gissning: grann} \\
            }
        }
    }
    & & \\ \\

    \theutterance \stepcounter{utterance}  
    & & & \multicolumn{2}{p{0.3\linewidth}}{
        \cellcolor[rgb]{0.9,0.9,0.9}{
            \makecell[{{p{\linewidth}}}]{
                \texttt{\tiny{[GM$|$GM]}}
                \texttt{attempts = 2} \\
\texttt{target\_word = panna} \\
\texttt{guess = grann} \\
\texttt{guess\_feedback = g<red> r<red> a<yellow> n<green> n<yellow>} \\
            }
        }
    }
    & & \\ \\

    \theutterance \stepcounter{utterance}  
    & & & \multicolumn{4}{p{0.6\linewidth}}{
        \cellcolor[rgb]{0.9,0.9,0.9}{
            \makecell[{{p{\linewidth}}}]{
                \texttt{\tiny{[P1$\langle$GM]}}
                \texttt{gissning\_återkoppling: g<röd> r<röd> a<gul> n<grön> n<gul>} \\
\\ 
\texttt{Ge ditt svar endast i detta format.} \\
\texttt{förklaring: detaljer} \\
\texttt{gissning: ord} \\
            }
        }
    }
     \\ \\

    \theutterance \stepcounter{utterance}  
    & \multicolumn{4}{p{0.6\linewidth}}{
        \cellcolor[rgb]{0.561,0.737,0.561}{
            \makecell[{{p{\linewidth}}}]{
                \texttt{\tiny{[P1$\rangle$GM]}}
                \texttt{förklaring: *g* och *r* är fel och utesluts. *a* finns med men på fel position. *n* är korrekt på fjärde platsen, och ett extra *n* finns någon annanstans (men inte på femte platsen, då det blev gult där). Nästa gissning testar nya bokstäver och placerar *a* och *n* på nya positioner.} \\
\texttt{gissning: kanna} \\
            }
        }
    }
    & & \\ \\

    \theutterance \stepcounter{utterance}  
    & & & \multicolumn{2}{p{0.3\linewidth}}{
        \cellcolor[rgb]{0.9,0.9,0.9}{
            \makecell[{{p{\linewidth}}}]{
                \texttt{\tiny{[GM$|$GM]}}
                \texttt{attempts = 3} \\
\texttt{target\_word = panna} \\
\texttt{guess = kanna} \\
\texttt{guess\_feedback = k<red> a<green> n<green> n<green> a<green>} \\
            }
        }
    }
    & & \\ \\

    \theutterance \stepcounter{utterance}  
    & & & \multicolumn{4}{p{0.6\linewidth}}{
        \cellcolor[rgb]{0.9,0.9,0.9}{
            \makecell[{{p{\linewidth}}}]{
                \texttt{\tiny{[P1$\langle$GM]}}
                \texttt{gissning\_återkoppling: k<röd> a<grön> n<grön> n<grön> a<grön>} \\
\\ 
\texttt{Ge ditt svar endast i detta format.} \\
\texttt{förklaring: detaljer} \\
\texttt{gissning: ord} \\
            }
        }
    }
     \\ \\

    \theutterance \stepcounter{utterance}  
    & \multicolumn{4}{p{0.6\linewidth}}{
        \cellcolor[rgb]{0.561,0.737,0.561}{
            \makecell[{{p{\linewidth}}}]{
                \texttt{\tiny{[P1$\rangle$GM]}}
                \texttt{förklaring: *k* är fel och utesluts. Målordet har *a* på andra platsen och *n* på tredje, fjärde och femte platsen. Strukturen är nu klar: `\_ a n n a`. Nästa gissning testar möjliga inledande bokstäver som bildar ett giltigt ord.} \\
\texttt{gissning: panna} \\
            }
        }
    }
    & & \\ \\

    \theutterance \stepcounter{utterance}  
    & & & \multicolumn{2}{p{0.3\linewidth}}{
        \cellcolor[rgb]{0.9,0.9,0.9}{
            \makecell[{{p{\linewidth}}}]{
                \texttt{\tiny{[GM$|$GM]}}
                \texttt{attempts = 4} \\
\texttt{target\_word = panna} \\
\texttt{guess = panna} \\
\texttt{guess\_feedback = p<green> a<green> n<green> n<green> a<green>} \\
            }
        }
    }
    & & \\ \\

    \theutterance \stepcounter{utterance}  
    & & & \multicolumn{2}{p{0.3\linewidth}}{
        \cellcolor[rgb]{0.9,0.9,0.9}{
            \makecell[{{p{\linewidth}}}]{
                \texttt{\tiny{[GM$|$GM]}}
                \texttt{game\_result = WIN} \\
            }
        }
    }
    & & \\ \\

\end{supertabular}
}

\end{document}
